% ---------------------------------------------------------------------------
% Beamer preamble for the model slides of the TP3 kick-off presentation.
% Geometry: aspectratio 169 (160mm x 90mm), rendered to PNG and placed
% full-bleed onto the 10in x 5.625in PowerPoint slides. Scale factor 1.5875.
% ---------------------------------------------------------------------------

% Arial-like sans serif text, matching the deck's Arial typography
\usepackage[scaled=0.95]{helvet}
\renewcommand{\familydefault}{\sfdefault}

% Classic LaTeX (Computer Modern) math, overriding beamer's sans-serif math
\usefonttheme[onlymath]{serif}

\usepackage{mathtools}
\usepackage{stmaryrd}
\usepackage{booktabs}
\usepackage{array}

% Iverson bracket and argmax, as in the report preamble
\DeclarePairedDelimiter{\ib}{\llbracket}{\rrbracket}
\DeclareMathOperator*{\argmax}{arg\,max}

% Plain white frames, no beamer furniture
\setbeamertemplate{navigation symbols}{}
\setbeamertemplate{footline}{}
\setbeamertemplate{headline}{}
\beamertemplatenavigationsymbolsempty

% Title: black, non-bold Arial, left aligned, matching the deck's 23pt titles
\setbeamercolor{frametitle}{fg=black,bg=white}
\setbeamerfont{frametitle}{size=\normalsize,series=\mdseries}
\setbeamertemplate{frametitle}{%
  \vspace{2.2mm}\hspace{-1mm}%
  \fontsize{14.5}{17}\selectfont\insertframetitle\par%
}

% Text block margins. The rendered slide is placed at 84% scale on the
% PowerPoint slide, so 6mm here lands the text close to the deck's own
% 0.83in title flush.
\setbeamersize{text margin left=6mm, text margin right=6mm}

% Component table: math symbol left, one-line explanation right.
% Rule-free, with generous leading between rows.
\newcommand{\componentcolsep}{8mm}
\newcommand{\componentrowsep}{2.4mm}
